\chapter{Het procesflowschema}
\label{h:pfd}

\section{Wat de installatie doet}
De installatie bestaat uit \VAR{I.n_units} parallelle ITM Neptune V-elektrolyse-units. Elke unit splitst
ultrazuiver water (UPW) met elektriciteit in waterstof en zuurstof. Rond die kern zitten vier deelsystemen:
\begin{enumerate}
  \item \textbf{\HH-lijn} (per unit): buffervat V-101, compressor K-101 van 20 naar 70\,barg, nakoeling in drie
        stappen (pakketkoeler, E-102, E-103), een kwaliteitscontrole, buffervat V-102 en drukregelaar PCV-101
        naar de \HH-backbone (66,2\,barg).
  \item \textbf{\OO-lijn} (per unit): de natte zuurstof (\OO\ met waterdamp) wordt in E-101 gekoeld tot
        30\degC. In knock-out-drum V-103 wordt het condensaat afgescheiden. Compressor K-102 brengt het gas op
        0,7\,barg naar de RWZI (rioolwaterzuivering), waar de zuurstof wordt gebruikt voor beluchting.
  \item \textbf{Water} (per unit): het condensaat uit V-103 gaat via pomp P-101 direct terug naar de unit (S16).
        Vers UPW (S1) vult alleen aan wat ontbreekt.
  \item \textbf{Warmteketen} (installatie): de restwarmte van de ITM-koelcircuits en de \HH-koelers gaat via één
        gesloten koelwatercircuit TW naar het warmtenet (direct via E-204 en via warmtepomp WP-201) en naar de
        industrie (via warmtepomp WP-202). Wat niet wordt afgenomen, gaat via de luchtkoelers AK naar buiten.
\end{enumerate}

\section{De tekeningen}
Figuur~\ref{f:pfd1} toont één unit en figuur~\ref{f:pfd2} de warmteketen. De getallen bij de stromen zijn de
resultaten van dit verslag (versie \VAR{"A" if A else "B"}).

\begin{landscape}
\begin{figure}[p]
\centering
\includegraphics[width=\linewidth,height=0.86\textheight,keepaspectratio]{fig/pfd_unit.pdf}
\caption{PFD blad 1: één ITM-unit. Alle debieten in kg/h per unit.}
\label{f:pfd1}
\end{figure}
\begin{figure}[p]
\centering
\includegraphics[width=\linewidth,height=0.86\textheight,keepaspectratio]{fig/pfd_warmte.pdf}
\caption{PFD blad 2: de warmteketen voor de hele installatie (8 units samen).}
\label{f:pfd2}
\end{figure}
\end{landscape}

\section{Apparatenlijst}
\begin{longtable}{p{2.1cm}p{4.2cm}p{9.0cm}}
\toprule
tag & naam & functie en waarom het er zit\\
\midrule\endhead
ITM & elektrolyser (×8) & splitst water in \HH\ en \OO; levert \HH\ op 30\,barg en \OO\ (nat) op 0,1\,barg\\
V-101 & lagedrukbuffer \HH & vangt het verschil op tussen de productie van de ITM en het debiet van de compressor; zorgt voor een stabiele zuigdruk (20–29\,barg)\\
K-101 & \HH-compressor & 20 $\rightarrow$ 70\,barg, zodat de backbone (66,2\,barg) bereikt wordt\\
PK-101 & pakketkoeler K-101 & koeler in het compressorpakket; koelt 30\,K af op het TW-circuit (warmteterugwinning)\\
E-102 & \HH-nakoeler & koelt \HH\ van $\approx$\,\VAR{n(R.T_S6,0)} naar 50\degC\ op TW (warmteterugwinning)\\
E-103 & \HH-eindkoeler & koelt \HH\ naar 25\degC\ op koud koelwater (TW is te warm daarvoor)\\
FV-101 & recycleklep & stuurt \HH\ terug naar V-101 als het debiet onder het minimum van K-101 komt (S9)\\
V-102 & hogedrukbuffer & dempt drukstoten naar de backbone; 30\,s levering\\
PCV-101 & drukregelaar & regelt terug van 70 naar 66,2\,barg (backbonedruk)\\
E-101 & \OO-voorkoeler & koelt natte \OO\ van 70 naar 30\degC; het grootste deel van de waterdamp condenseert\\
V-103 & KO-drum & scheidt het condensaat van het gas\\
K-102 & \OO-compressor & 0,1 $\rightarrow$ 0,7\,barg voor de beluchting op de RWZI\\
PK-102 & pakketkoeler K-102 & koelt 30\,K na de compressie\\
P-101 & condensaatpomp & pompt het condensaat (UPW-kwaliteit) terug naar de unit (S16)\\
E-201 & platenwisselaar (×8) & haalt warmte uit het gesloten ITM-koelcircuit naar TW; het ITM-circuit blijft gesloten\\
AK & luchtkoeler (×8) & bestaande koeler van het ITM-circuit; trim en back-up als TW niet alles afneemt\\
E-204 & directe voorverwarming & TW verwarmt het netwater van 40 naar 45\degC, zonder arbeid\\
E-202 / E-206 & verdampers & WP-201 en WP-202 halen warmte uit TW\\
K-201 / K-202 & warmtepompcompressoren & brengen het koudemiddel op condensatiedruk\\
E-203 / E-205 & condensors & leveren warmte aan het net (70\degC) en aan de industrie (95\degC)\\
TV-201 / TV-202 & smoorventielen & expansie van condensatie- naar verdampingsdruk ($h$ constant)\\
\bottomrule
\end{longtable}

\section{Stroomtabel}
\label{s:stroomtabel}
Debiet, temperatuur en druk van alle stromen. ``Unit'' betekent per unit; ``inst.'' betekent voor de hele installatie.
\begin{center}\footnotesize\setlength{\tabcolsep}{3.5pt}
\begin{longtable}{llp{5.2cm}p{2.3cm}rrr}
\toprule
stroom & schaal & omschrijving & medium & $\md$ [kg/h] & $T$ [\textdegree C] & $P$ [kPa] (barg)\\
\midrule\endhead
%% for L in R.leidingen
\VAR{L.id} & \VAR{"unit" if L.schaal == "per unit" else "inst."} & \VAR{tex(L.naam)} & \VAR{tex(L.medium)} & \VAR{n(L.m, 1 if L.m < 1000 else 0)} & \VAR{n(L.T, 1)} & \VAR{(n(L.p*100+PATM,1) ~ " (" ~ n(L.p,1) ~ ")") if L.p is not none else "--"}\\
%% endfor
\bottomrule
\end{longtable}
\end{center}
Bij vollast staat recyclestroom S9 dicht (0\kgh). In de tabel staat het maximale ontwerpdebiet, dat nodig is
voor de leidingdimensionering. S10 (blowoff) is een veiligheidsafblaas en voert normaal niets af. Bij TW, N en I
is de druk niet in het model opgenomen ($h$ van vloeistof hangt nauwelijks van $P$ af).
